L’apprentissage automatique (Machine Learning, ML) a connu une progression fulgu-
rante, menant `a des mod`eles d’une performance remarquable, souvent sup´erieure `a celle
de l’humain dans des tˆaches sp´ecifiques. Les r´eseaux de neurones profonds (Deep Neural
Networks, DNN), en particulier, ont r´evolutionn´e des domaines comme la vision par or-
dinateur, le traitement du langage naturel et la reconnaissance vocale. Cependant, cette
performance s’accompagne d’un coˆut significatif : une complexit´e interne croissante qui
rend leurs m´ecanismes de d´ecision extrˆemement difficiles, voire impossibles, `a comprendre
pour un observateur humain. Ces mod`eles sont fr´equemment qualifi´es de ”boˆıtes noires”
(black boxes), des syst`emes dont les rouages internes, qui transforment une entr´ee en une
sortie, sont opaques et inaccessibles `a l’analyse directe [1, 2].


Cette opacit´e, bien que tol´erable dans des applications `a faible enjeu comme la recom-
mandation de publicit´es en ligne, devient une pr´eoccupation majeure lorsque ces syst`emes
sont d´eploy´es dans des domaines sensibles o`u les d´ecisions ont des cons´equences financi`eres,
l´egales, ou vitales. Des secteurs comme la sant´e, la finance, la justice p´enale ou les v´ehicules
autonomes exigent un niveau ´elev´e de confiance, de responsabilit´e et de transparence [1, 2].

C’est dans ce contexte qu’a ´emerg´e le domaine de l’Intelligence Artificielle Explicable
(Explainable AI, XAI). L’objectif de l’XAI n’est pas de sacrifier la performance au profit de
la simplicit´e, mais de d´evelopper des m´ethodes et des techniques permettant de rendre les
d´ecisions des mod`eles complexes plus transparentes, compr´ehensibles et justifiables pour
les humains [1]. La demande pour l’XAI n’est donc pas une simple curiosit´e acad´emique ;
elle est une cons´equence directe du succ`es et de l’int´egration croissante de l’IA dans des
syst`emes socio-techniques critiques. Plus l’impact et l’autonomie de l’IA augmentent, plus
son opacit´e devient intol´erable. Le probl`eme n’est pas tant que le mod`ele est une boˆıte
noire, mais que la soci´et´e, `a juste titre, lui demande des comptes pour des d´ecisions qui,
par nature, exigent une justification.


Confiance et Adoption : Les utilisateurs, en parti-
culier dans les domaines critiques, hésitent à faire
confiance à des décisions qu’ils ne comprennent pas.
— Débogage et Amélioration : Sans compréhension
du fonctionnement interne, il est difficile d’identifier
et de corriger les erreurs du modèle.
— Équité et Biais : Les modèles peuvent apprendre
et perpétuer des biais présents dans les données,
conduisant à des décisions discriminatoires. L’inter-
prétabilité est essentielle pour détecter et atténuer ces
biais.
— Responsabilité et Réglementation : En cas d’er-
reur du modèle ayant des conséquences graves, il est
crucial de pouvoir déterminer la responsabilité. Des
réglementations comme le RGPD en Europe intro-
duisent un "droit à l’explication".

 Découverte Scientifique : Dans des domaines
comme la biologie ou la médecine, l’interprétabilité
peut aider à découvrir de nouvelles connaissances à
partir des patterns identifiés par le modèle.
La "découverte" dans ce contexte se réfère à la prise de
conscience progressive de ces enjeux et à l’émergence de
l’interprétabilité (ou XAI - Explainable Artificial Intelli-
gence) comme un champ de recherche majeur.


analyse les principales méthodes d’interprétabilité,
leur fonctionnement, leurs applications, forces et limita-
tions, en s’appuyant notamment sur les travaux de Doran
et al. (2017) et Burkart Huber (2021). La seconde par-
tie introduit les systèmes de recommandation, leurs tech-
niques fondamentales, et propose une comparaison dé-
taillée de deux approches majeures, en se référant aux sur-
veys comme celui de Ko et al. et aux avancées en filtrage
collaboratif décrites par Koren Bell.


Doran et al. (2017) proposent une conceptualisation de
l’explicabilité en IA qui distingue différentes perspectives.
Ils soulignent que ce qui constitue une "bonne" explication
peut varier considérablement en fonction de l’audience
(par exemple, un développeur de modèles vs un utilisateur
final non expert) et du but de l’explication. Ils identifient
des facettes telles que la transparence du modèle, la justi-
fication post-hoc, et la simulation du raisonnement.
Burkart and Huber (2021) offrent un survey complet sur
l’explicabilité des modèles d’apprentissage supervisé. Ils
classifient les méthodes d’explicabilité selon plusieurs di-
mensions :
— Intrinsèque vs. Post-hoc : L’interprétabilité intrin-
sèque est obtenue avec des modèles simples par
conception (ex : arbres de décision, régression li-
néaire). Les méthodes post-hoc sont appliquées à des
modèles déjà entraînés, souvent complexes.
— Spécifique au modèle vs. Agnostique au modèle :
Les méthodes spécifiques ne fonctionnent qu’avec
certains types de modèles (ex : attention pour les
réseaux de neurones). Les méthodes agnostiques
peuvent être appliquées à n’importe quel modèle (ex :
LIME, SHAP).
— Locale vs. Globale : Les explications locales se
concentrent sur une prédiction individuelle. Les ex-
plications globales visent à comprendre le comporte-
ment général du modèle.
Ils discutent également des différentes formes que peuvent
prendre les explications (règles, caractéristiques impor-
tantes, exemples prototypiques, etc.).